\documentclass[10pt,a4paper]{amsart}
\usepackage[margin=19mm]{geometry}
\usepackage{amsmath,amssymb,mathtools}
\usepackage[scr=rsfs,cal=euler]{mathalfa}
\usepackage{xfrac}
\usepackage[unicode,hidelinks,bookmarks=false]{hyperref}
\newcommand{\ie}{i.e.\ }
\newcommand{\cf}{cf.\ }
\newcommand{\resp}{respectively\ }
\newcommand{\Subsection}{\subsection}
\providecommand{\nobreakdash}{\nobreak\hbox{-}}
\setlength{\emergencystretch}{2em}
\multlinegap0pt
\usepackage{bm}
\usepackage{bbm}
\usepackage{dsfont}
\usepackage{xspace}
\usepackage{cancel}
\usepackage{mathtools}
\usepackage{amsthm}
\makeatletter
\@ifundefined{thm}{\newtheorem{thm}{Thm}}{}
\@ifundefined{prop}{\newtheorem{prop}[thm]{Proposition}}{}
\makeatother
\title[Periodicity of point processes in abelian groups without lat]{Periodicity of point processes in abelian groups without lattices}
\author{Nachi Avraham-Re'em, Michael Bj\"orklund,, Rickard Cullman}
\date{Source: arXiv:2509.20847v2 (CC BY 4.0)}
\begin{document}
\maketitle

\noindent\emph{Source and licence.} Periodicity of point processes in abelian groups without lattices. Version arXiv:2509.20847v2 (CC BY 4.0), distributed under the Creative Commons Attribution 4.0 International licence, as recorded in the arXiv licence metadata for this exact version and shown on the abstract page https://arxiv.org/abs/2509.20847v2. Full rights notice: licenses/arxiv-2509.20847-CC-BY-4.0.txt.\par\smallskip

\noindent\textit{Editorial scope.} Counted unit: the Prop whose source label is \texttt{None} in the article (source0.tex lines 1368--1377, source label \texttt{None}), delivered together with the article's own complete original proof of it (source0.tex lines 1379--1405, one \texttt{proof} environment of 27 lines). No mathematics is written, repaired or re-derived: statement and proof are the article's own text line for line. Required context retained uncounted: none. Every definition, display and label of those lines is retained unchanged. Omitted from this package and not redistributed: 1--1367, 1378--1378, 1406--1411. Nothing inside a retained excerpt is omitted or altered: every retained body is delivered exactly as the article's lines are written, so no omission receipt of any kind accompanies this package. Citations: the retained excerpts carry no citation command at all, so no bibliography entry of the article is transcribed and no reference is redistributed. Reference adaptation: none was needed and none was made; every reference inside the delivered unit resolves inside it, so no unresolved reference or citation is produced. Printed slips kept literal: no printed word, symbol, hypothesis or number of the counted statement or of its proof was altered. Layout and class: the article's own class is \texttt{amsart} with packages including geometry, amssymb, url, amsmath, amsthm, mathrsfs, appendix, inputenc, babel, cite, hyperref, cleveref, placeins, textgreek; the wrapper is the lane's pinned \texttt{amsart} editorial wrapper at 10pt on A4, which restates the theorem-like environments the retained text needs and adds the article's own macro definitions that the retained text needs (none), with extra wrapper packages (bm, bbm, dsfont, xspace, cancel, mathtools). The delivered numbering is the unit's own. Assets: no retained span contains a figure environment, a graphics-inclusion command, a PDF-page inclusion, a table or a TikZ picture; no image file is copied into this package, the package declares no asset dependency and no image file is redistributed with it. Licence: version arXiv:2509.20847v2 of this article is distributed under the Creative Commons Attribution 4.0 International licence, as recorded in the arXiv licence metadata for this exact version and shown on the abstract page https://arxiv.org/abs/2509.20847v2. A full rights notice is delivered at licenses/arxiv-2509.20847-CC-BY-4.0.txt. Characters: every retained excerpt is pure ASCII, so the delivered engine, XeLaTeX with its default Latin Modern text font, reports no missing character.
\begin{prop}
Let \(G\) be a \hyperlink{classQ}{class \(\mathcal{Q}\)} group satisfying the following properties:
\begin{enumerate}
    \item \(G\) is torsion-free and divisible.
    \item There is a compact open subgroup \(U<G\) with no nontrivial finite subgroups, and \(U\backslash G\) is torsion.
\end{enumerate}
Let \(Q\) be the set of orders of elements in \(U\backslash G\). Then every lattice \(\Gamma< G\times\mathbb{R}\) is of the form
\[\Gamma=\mathbb{Z}\big[1/Q\big]\left(u,t\right)=\big\{ \big({\textstyle \frac{m}{q}}u,{\textstyle\frac{m}{q}}t\big):m\in\mathbb{Z},q\in Q\big\}<G\times\mathbb{R},\footnote{The writing \(g/q\) for \(g\in G\) and \(q\in\mathbb{Z}\backslash\{0\}\) is justified since \(G\) is divisible and torsion free.}\]
for some \(\left(u,t\right)\in U\times\mathbb{R}\).
\end{prop}

\begin{proof}
Let \(\Gamma<G\times\mathbb{R}\) be a lattice and put \(\Gamma_{U}\coloneqq\Gamma\cap\left(U\times\mathbb{R}\right)\). We claim that the projection map \(\operatorname{proj}_{\mathbb{R}}:\Gamma_{U}\to\operatorname{proj}_{\mathbb{R}}\left(\Gamma_{U}\right)\) forms an isomorphism of \(\Gamma_{U}\) onto a discrete subgroup of \(\mathbb{R}\). Indeed, \(\ker\operatorname{proj}_{\mathbb{R}}\mid_{\Gamma_{U}}=\Gamma\cap\left(U\times\left\{ 0\right\} \right)\) is a finite subgroup of \(U\) (since \(\Gamma\) is discrete and \(U\) is compact), and since \(U\) admits no nontrivial finite subgroups this kernel is trivial. We have left to show that \(\mathrm{proj}_{\mathbb{R}}\left(\Gamma_{U}\right)\) is a discrete subgroup of \(\mathbb{R}\). Suppose \(\left(t_{n}\right)_{n\in\mathbb{N}}\subset\mathrm{proj}_{\mathbb{R}}\left(\Gamma_{U}\right)\) is a sequence with \(t_{n}\to0\) as \(n\to\infty\). Then there exists a sequence \(\left(u_{n}\right)_{n\in\mathbb{N}}\subset U\) such that \(\gamma_{n}\coloneqq\left(u_{n},t_{n}\right)\in\Gamma\) for every \(n\in\mathbb{N}\), and since \(U\) is compact, we may assume (after passing to a subsequence) that \(u_{n}\to u\) as \(n\to\infty\) for some \(u\in U\), and consequently \(\gamma_{n}\to\left(u,0\right)\) as \(n\to\infty\). Because \(\Gamma\) is discrete and \(\left(\gamma_{n}\right)_{n\in\mathbb{N}}\subset\Gamma\), it follows that \(\gamma_{n}=\left(u,0\right)\) for all sufficiently large \(n\), so \(t_{n}=0\) for all sufficiently large \(n\). Thus every sequence in \(\mathrm{proj}_{\mathbb{R}}\left(\Gamma_{U}\right)\) that converges to \(0\) is eventually constant, hence \(\mathrm{proj}_{\mathbb{R}}\left(\Gamma_{U}\right)\) is a discrete subgroup of \(\mathbb{R}\).

We conclude that \(\mathrm{proj}_{\mathbb{R}}\mid_{\Gamma_{U}}:\Gamma_{U}\to\mathrm{proj}_{\mathbb{R}}\left(\Gamma_{U}\right)\) is an isomorphism of discrete groups. Hence, by the structure of discrete subgroups of \(\mathbb{R}\), there is \(\left(u,t\right)\in U\times\mathbb{R}\) such that
\[
\Gamma_{U}=\mathbb{Z}\left(u,t\right)=\bigl\{(mu,mt):m\in\mathbb{Z}\bigr\}.
\]
We now claim that the following identity holds:
\begin{equation}\label{eq:gammaZQ}
\Gamma=\mathbb{Z}\big[1/Q^{\prime}\big]\left(u,t\right),
\end{equation}
where \(Q^{\prime}\) denotes the set of orders of elements in \(\left\{Ug:g\in\mathrm{proj}_{G}\left(\Gamma\right)\right\}\subseteq U\backslash G\). Let \((\gamma_{1},\gamma_{2})\in\Gamma\) be arbitrary. Because \(U\backslash G\) is torsion, there is \(q\in Q^{\prime}\) with \(q\gamma_{1}\in U\) (note \(\gamma_{1}\in\mathrm{proj}_{G}\left(\Gamma\right)\)), and then
\[\left(q\gamma_{1},q\gamma_{2}\right)\in\Gamma_{U}=\mathbb{Z}\left(u,t\right),\]
so \(\left(q\gamma_{1},q\gamma_{2}\right)=\left(mu,mt\right)\) for some \(m\in\mathbb{Z}\), and using the divisibility of \(G\) we get
\[\left(\gamma_{1},\gamma_{2}\right)=\frac{m}{q}\left(u,t\right)\in\mathbb{Z}\big[1/q\big]\left(u,t\right)\subseteq\mathbb{Z}\big[1/Q^{\prime}\big]\left(u,t\right).\]
Conversely, let \(q\in Q^{\prime}\) be arbitrary and pick \(g\in\mathrm{proj}_{G}\left(\Gamma\right)\) such that \(Ug\in U\backslash G\) has order \(q\), thus \(qg\in U\). Pick some \(s\in\mathbb{R}\) with \(\left(g,s\right)\in\Gamma\), and as above we may write
\[\left(qg,qs\right)=m\left(u,t\right)\quad\text{for some }m\in\mathbb{Z},\]
so \(qg=mu\in U\). The minimality of \(q\) forces \(\gcd\left(m,q\right)=1\) (if \(d\coloneqq\gcd\left(m,q\right)>1\) then from \(qg=mu\) we get \(\left(q/d\right)g=\left(m/d\right)u\in U\), contradicting the minimality of \(q\)). Then pick \(a,b\in\mathbb{Z}\) with \(am+bq=1\) (B\'{e}zout's identity), and using the divisibility of \(G\) we get
\[\frac{1}{q}\left(u,t\right)=a\left(g,s\right)+b\left(u,t\right)\in\Gamma,\]
so each generator \(\frac{1}{q}\left(u,t\right)\in\mathbb{Z}\left[1/Q^{\prime}\right]\left(u,t\right)\) lies in \(\Gamma\). This establishes \eqref{eq:gammaZQ}.

Finally, let \(Q\) be the set of orders of elements in \(U\backslash G\), and we will prove that \(Q=Q^{\prime}\). The inclusion \(Q^{\prime}\subseteq Q\) is clear, and for the reverse inclusion it suffices to show that \(U\mathrm{proj}_{G}\left(\Gamma\right)=G\). Since \(\Gamma\) is a lattice in the abelian group \(G\times\mathbb{R}\), there is a compact set \(C\subset G\times\mathbb{R}\) with \(\Gamma C=G\times\mathbb{R}\), hence
\[G=\mathrm{proj}_{G}\left(\Gamma\right)K,\text{ where }K\coloneqq\mathrm{proj}_{G}\left(C\right).\]
Since \(U\) is compact open and \(K\) is compact, there is a finite set \(F\subseteq K\) such that \(UK=UF\), hence
\[G=\mathrm{proj}_{G}\left(\Gamma\right)K\subseteq\left(U\mathrm{proj}_{G}\left(\Gamma\right)\right)\left(UK\right)=\left(U\mathrm{proj}_{G}\left(\Gamma\right)\right)F,\]
which means that \(U\mathrm{proj}_{G}\left(\Gamma\right)\) is a finite index subgroup of \(G\). Since \(G\) is torsion-free and divisible, it has no proper finite index subgroups, and therefore \(U\mathrm{proj}_{G}\left(\Gamma\right)=G\). This completes the proof.
\end{proof}

\end{document}
